\begin{textsample}{POS Dim 3 – sc – Score 9.00 – sc/sc\_em\_12.txt}  \label{ex:f3_pos_138}
4.2.4 Organizador curricular do componente Projeto de Vida Alinhado às diretrizes da BNCC, o componente Projeto de Vida para o Território Catarinense tem por foco o desenvolvimento de habilidades e a promoção de competências, a se alcançar a partir da elaboração conceitual, mobilizada pelos objetos de conhecimento.
Neste contexto, o organizador curricular a ser apresentado nas próximas páginas está alicerçado nas dez competências gerais da BNCC.
Esta organização é propositiva, podendo ser usada de forma flexível, com o objetivo de promover percursos de aprendizagem que contribuam para a construção dos projetos de vida dos estudantes.
A construção deste organizador está pautada no preceito da aprendizagem progressiva e espiralada, propondo -se a elaboração conceitual a partir da complexificação dos objetos do conhecimento para o desenvolvimento de habilidades e competências.
É válido destacar que este organizador pode ser usado de diversas formas, orientando a criação de sequências didáticas, de acordo com as habilidades ou objetivos sobre os quais se deseja trabalhar.
Quadro 10 - Matriz curricular para o componente Projeto de Vida 10 competências gerais da BNCC 1.
Conhecimento Competência geral : Valorizar e utilizar os conhecimentos historicamente construídos sobre o mundo físico, social, cultural e digital para entender e explicar a realidade, continuar aprendendo e colaborar para a construção de uma sociedade justa, democrática e inclusiva.
Articulação com o projeto de vida : Esta competência se articula ao projeto de vida mediante a necessidade dos estudantes de conhecerem seus contextos e as oportunidades neles inscritas, de modo a fazer escolhas orientadas pelo conhecimento que têm do mundo físico, social, cultural e digital.
Habilidades : - Identificar e valorizar os conhecimentos, as habilidades e as competências desenvolvidos no percurso formativo. - Compreender a influência do contexto socioeconômico, cultural, histórico e político e aplicá -la ao próprio contexto pessoal. - Identificar os aspectos que conferem sentido à vida humana e à própria existência, em particular. - Compreender a estrutura dos projetos de vida, diferenciando -os de sonhos e fantasias. - Reconhecer, identificar e valorizar a história e a cultura local na construção de seu projeto de vida. - Problematizar o enfrentamento dos desafios dos jovens e sensibilizar para o fazer. - Promover ações de valorização do aprendizado e da vivência escolar.
Objetos de conhecimento : - Trajetória escolar : vivências escolares, aprendizados e atores envolvidos ( professores, \textbf{colegas}, familiares e demais membros da comunidade ) que constituem a formação integral. - Influência dos diversos contextos — socioeconômico, cultural, histórico e político — na construção do projeto de vida. - Sentidos da vida : aspectos que conferem sentido à vida e promovem realização pessoal, felicidade, superação e enfrentamento de condições adversas. - Diferença entre projetos de vida ( aspirações articuladas a planejamento e engajamento ), sonhos ( aspirações desarticuladas de planejamento e engajamento ) e fantasias ( \textbf{desejos} desconectados da realidade ). - História e cultura local e suas influências na construção dos projetos de vida dos estudantes catarinenses. - Diversidade e desafios das juventudes. 2.
Pensamento científico, crítico e criativo Competência geral : Exercitar a \textbf{curiosidade} intelectual e recorrer à abordagem própria das ciências, incluindo a investigação, a reflexão, a análise crítica, a imaginação e a criatividade, para investigar causas, elaborar e testar hipóteses, formular e resolver problemas e criar soluções ( inclusive tecnológicas ) com base nos conhecimentos das diferentes áreas.
Articulação com o projeto de vida : Esta competência se articula com o projeto de vida mediante a necessidade dos estudantes de construir trajetórias motivadas pelo senso crítico, pela reflexão, pela investigação das circunstâncias que os afetam e pela proposição criativa de trajetórias que visem à conquista de suas aspirações.
Habilidades : - Investigar e analisar dados e informações, formular hipóteses e relacionar variáveis na identificação de problemas sociais e ambientais que possam motivar a construção de seus projetos de vida. - Usar a criatividade para criar soluções inovadoras para demandas sociais e ambientais no entorno escolar e, quando possível, aplicá -las. - Usar a criatividade para elaborar um planejamento viável, desejado e efetivo para a realização do seu projeto de vida. - Investigar diversas possibilidades de escolha profissional para contribuir com a tomada de decisão responsável. - Analisar as funções sociais das profissões e suas relações com as demandas da sociedade contemporânea e do mercado de trabalho. - Promover intervenções na escola, ou em seu entorno, que evidenciem o protagonismo juvenil, por meio do desenvolvimento de projetos de interesse dos estudantes, que contribuam com seus projetos de vida.
Objetos de conhecimento : - Problemas sociais e ambientais locais e globais. - Justiça social. - Democracia. - Sustentabilidade ambiental. - Inovação como prática de desenvolvimento humano na dimensão pessoal, cidadã e profissional. - Carreiras profissionais : função social, formação, remuneração, reconhecimento social, competências técnicas e comportamentais. - Mercado de trabalho e demandas em diferentes escalas ( local, regional, nacional e internacional ). - Critérios para um planejamento eficaz : viabilidade, desejabilidade e efetividade. 3.
Repertório cultural Competência geral : Valorizar as diversas manifestações artísticas e culturais, das locais às mundiais, e delas fruir e também participar de práticas diversificadas da produção artístico-cultural.
Articulação com o projeto de vida : Esta competência se articula com o projeto de vida mediante a necessidade dos estudantes de reconhecerem a influência da cultura na qual estão inseridos em seus projetos de vida, \textbf{apoiando} -se nela para se inspirarem a construir formas satisfatórias de viver a vida.
Habilidades : - Reconhecer, respeitar e valorizar manifestações culturais, pessoas e ideias de diferentes regionalidades e territorialidades. - Sistematizar o projeto de vida utilizando diferentes recursos artísticos. - Promover ações que reconheçam a diversidade cultural catarinense, bem como suas manifestações artísticas. - Conhecer e aplicar os princípios da comunicação não violenta nas relações interpessoais. - Conhecer e aplicar os princípios da cultura de paz para dialogar com pessoas que pensam de forma diferente em seus locais de vivência.
Objetos de conhecimento : - Manifestações culturais e artísticas de diversas escalas ( local, regional, nacional e internacional ) que podem influenciar a construção dos projetos de vida. 4.
Comunicação Competência geral : Utilizar diferentes linguagens – verbal ( oral ou visual-motora, como Libras, e escrita ), corporal, visual, sonora e digital –, bem como conhecimentos das diversas linguagens ( artística, matemática e científica ), para se expressar e partilhar informações, experiências, ideias e sentimentos em diferentes contextos e produzir sentidos que levem ao entendimento mútuo.
Articulação com o projeto de vida : Esta competência se articula com o projeto de vida mediante a necessidade dos estudantes de aprenderem a comunicar suas ideias, seus pensamentos e sentimentos de forma respeitosa e empática, o que fortalece a \textbf{autoconfiança} em suas decisões e angaria o apoio de outras pessoas para a concretização de seus objetivos.
Habilidades : - Utilizar estratégias de resolução de conflitos em diversos contextos. - Conhecer e aplicar estratégias da \textbf{escuta} empática. - Expressar os próprios pensamentos e sentimentos com clareza e segurança. - Propor ações que desenvolvam o aperfeiçoamento da comunicação.
Objetos de conhecimento : - Princípios da comunicação não violenta. - Cultura de paz. - Resolução de conflitos. - \textbf{Escuta} empática. 5.
Cultura digital Competência geral : Compreender, utilizar e criar tecnologias digitais de informação e comunicação de forma crítica, significativa, reflexiva e ética nas diversas práticas sociais ( incluindo as escolares ) para se comunicar, acessar e disseminar informações, produzir conhecimentos, resolver problemas e exercer protagonismo e autoria na vida pessoal e coletiva.
Articulação com o projeto de vida : Esta competência se articula com o projeto de vida mediante a necessidade dos estudantes de se inserirem eticamente na sociedade contemporânea, cujas relações e processos estão cada vez mais marcados pelas tecnologias digitais, as quais, inclusive, configuram um importante mercado de trabalho no qual se podem inserir.
Habilidades : - Aplicar princípios éticos em suas interações nos ambientes digitais. - Praticar o autocontrole para lidar com o excesso de tempo dedicado às redes sociais. - Utilizar a ferramenta do podcast para aprender sobre temas de interesse pessoal e divulgar ideias. - Usar as redes sociais como plataforma para a defesa de causas sociais. - Compreender os riscos envolvidos na exposição e compartilhamento de dados e informações nos meios digitais. - Compreender os fatores envolvidos no cyberbullying e combatê -los ativamente. - Refletir sobre o comportamento humano nas redes sociais. - Utilizar redes sociais para explorar o mundo do trabalho e o universo das profissões. - Promover ações de uso consciente das redes em ambiente escolar. - Conhecer a legislação do ambiente virtual, problematizar e conscientizar para problemas como as fake news, deep web, entre outros.
Objetos de conhecimento : - Responsabilidade nas redes sociais. - Engajamento social e divulgação de ideias nas redes digitais. - Uso de ferramentas de busca e pesquisa para o aprendizado de temas de interesse. - Divulgação de dados e informações : riscos. - Cyberbullying. - Criação e compartilhamento de fake news : riscos. - Comportamentos nas redes sociais digitais : cancelamento, stalker. - Redes sociais profissionais : LinkedIn e ResearchGate. 6.
Trabalho e projeto de vida Competência geral : Valorizar a diversidade de saberes e vivências culturais e apropriar -se de conhecimentos e experiências que lhe possibilitem entender as relações próprias do mundo do trabalho e fazer escolhas alinhadas ao exercício da cidadania e a seu projeto de vida, com liberdade, autonomia, consciência crítica e responsabilidade.
Articulação com o projeto de vida : Esta competência se articula com o projeto de vida mediante a necessidade dos estudantes de compreenderem a si mesmos e o espaço que eles ocupam no mundo social, bem como o conjunto de saberes e experiências que podem contribuir para que criem projetos de vida com liberdade, autonomia, criticidade e responsabilidade.
Habilidades : - Reconhecer as relações entre os itinerários formativos e o mundo do trabalho. - Compreender que a formação escolar contribui para o desenvolvimento pessoal e permite a conquista de objetivos profissionais. - Conhecer e desenvolver procedimentos de planejamento, execução e \textbf{acompanhamento} de ações para o cumprimento de objetivos e metas de curto, médio e longo prazo, no âmbito pessoal e coletivo. - Reconhecer os fatores pessoais e sociais que interferem na tomada de decisões e fazer escolhas que impactam a vida pessoal e/ou coletiva de forma autônoma, \textbf{criteriosa} e ética. - Identificar caminhos e estratégias para superar as dificuldades e alicerçar a busca da realização de seus objetivos, com determinação e resiliência. - Identificar as características, as possibilidades e os desafios do trabalho no século XXI. - Relacionar interesses, habilidades, conhecimentos e oportunidades que correspondem às aspirações futuras com a inserção no mercado de trabalho. - Reconhecer e valorizar o papel da família no desenvolvimento físico, cognitivo, afetivo e social e a importância de seu apoio na construção do projeto de vida. - Compreender que redes de apoio auxiliam na superação de dificuldades para a realização de seus projetos e articular contatos com estas possíveis redes.
Objetos de conhecimento : - Relações entre os itinerários formativos e o mundo do trabalho. - Importância da formação para a construção de uma carreira profissional. - Diferenças entre objetivos, metas e estratégias. - Planejamento e autogestão. - Escolhas individuais e coletivas. - Critérios para a escolha profissional. - Trabalho no mundo contemporâneo. - \textbf{Equilíbrio} entre vida profissional e vida pessoal. - Objetivos de crescimento pessoal. - Potencial econômico e social da comunidade e da região. - Função social da família e apoio ao projeto de vida. - Redes de apoio para a construção do projeto de vida. - Modalidades de ensino e ofertas de formação em nível técnico e superior. 7.
Argumentação Competência geral : Argumentar com base em fatos, dados e informações confiáveis, para formular, negociar e defender ideias, pontos de vista e decisões comuns que respeitem e promovam os direitos humanos, a consciência socioambiental e o consumo responsável em âmbito local, regional e global, com posicionamento ético em relação ao \textbf{cuidado} de si mesmo, dos outros e do planeta.
Articulação com o projeto de vida : Esta competência se articula com o projeto de vida mediante a necessidade dos estudantes de argumentarem sobre seus pontos de vista em relação às decisões que fundamentam seus projetos de vida ao entrar em contato com ideias diferentes das suas, de modo a podê -las defender e a aprender com a perspectiva alheia, de modo ético e responsável.
Habilidades : - Argumentar, com base em fatos e dados, para negociar e defender ideias e pontos de vista em situação de debate público e privado. - Exercitar e praticar a tomada de perspectiva e a empatia para reconhecer e compreender as ideias, pensamentos, sentimentos e comportamentos alheios. - Utilizar linguagem verbal e corporal respeitosa ao argumentar sobre seus pontos de vista. - Argumentar sobre as próprias escolhas, explicando quais critérios foram utilizados para fazê -las.
Objetos de conhecimento : - Características e diferenças entre argumentação e persuasão. - Veracidade de dados e informações. - Tomada de perspectiva e empatia. - Defesa das próprias escolhas. 8.
Autoconhecimento e autocuidado Competência geral : Conhecer -se, apreciar -se e cuidar de sua saúde física e emocional, compreendendo -se na diversidade humana e reconhecendo suas emoções e as dos outros, com autocrítica e capacidade para lidar com elas.
Articulação com o projeto de vida : Esta competência se articula com o projeto de vida mediante a necessidade dos estudantes de conhecerem a si mesmos para poderem tomar decisões responsáveis sobre quem desejam ser no futuro.
Estas decisões também devem ser pautadas no autocuidado, que é a capacidade de não se expor a riscos nas esferas da saúde física e da saúde mental.
Habilidades : - Identificar atributos da própria personalidade. - Reconhecer valores, pensamentos, sentimentos e \textbf{hábitos} e regular as próprias condutas. - Compreender que a autoestima é influenciada por fatores subjetivos e socioculturais e ser capaz de elaborar uma representação positiva de si mesmo. - Reconhecer e avaliar as características que constituem a própria identidade, relacionando -as com o seu projeto de vida. - Compreender que a identidade é uma construção que cada indivíduo realiza ao longo de sua vida na interação com o meio, sob a influência de elementos políticos e socioculturais, entre outros. - Elaborar a narrativa autobiográfica, atribuindo significados às experiências de vida a partir de uma análise consciente de seu papel na construção da identidade e do projeto de vida. - Promover ações de \textbf{cuidado} com a saúde física e emocional no ambiente escolar.
Objetos de conhecimento : - Conceito de personalidade. - Conceito e importância da autoestima. - Identidade pessoal e identidade social. - Emoções e sentimentos agradáveis e desagradáveis. - Autorregulação. - Valores pessoais e desejáveis de universalização. - Narrativa de vida. - Práticas de autocuidado. 9.
Empatia e cooperação Competência geral : Exercitar a empatia, o diálogo, a resolução de conflitos e a cooperação, fazendo -se respeitar e promovendo o respeito ao outro e aos Direitos Humanos, com acolhimento e valorização da diversidade de indivíduos e de grupos sociais, seus saberes, identidades, culturas e potencialidades, sem preconceitos de qualquer natureza.
Articulação com o projeto de vida : Esta competência se articula com o projeto de vida mediante a necessidade dos estudantes de compreenderem que, mesmo que o projeto de vida seja uma construção individual, ele se realiza na esfera social, na relação com outras pessoas e depende do apoio e da cooperação de múltiplos agentes.
Além disso, a capacidade de empatizar e cooperar com outras pessoas pode levar à criação de projetos de vida coletivos, voltados ao compromisso com o social.
Habilidades : - Reconhecer a importância do convívio com o outro. - Promover ações de enfrentamento de preconceito, discriminação, intolerância e as diferentes formas de violências. - Praticar a tomada de perspectiva e a empatia para reconhecer e compreender ideias, pensamentos, sentimentos e comportamentos alheios. - Agir com empatia, sendo capaz de assumir a perspectiva dos outros, compreendendo as necessidades e sentimentos alheios, construindo relacionamentos baseados no compartilhamento e na abertura para o convívio. - Promover ações que valorizem o direito à diferença e reconhecer as identidades, valorizando o multiculturalismo e a inclusão. - Compreender o protagonismo e a identidade dos diversos sujeitos históricos ( \textbf{crianças}, juventudes, idosos, mulheres, população LGBTQIA+, negros, quilombolas, indígenas, pessoas com deficiência, povos originários, entre outros ).
Objetos de conhecimento : - Diversidade de modos de viver a vida. - Preconceitos e intolerâncias. - Relações competitivas e cooperativas. - Protagonismo juvenil. - Conceito de alteridade. - Lutas e pautas identitárias ( \textbf{crianças}, juventudes, idosos, mulheres, população LGBTQIA+, negros, quilombolas, indígenas, pessoas com deficiência, povos originários, entre outros ). 10.
Responsabilidade e cidadania Competência geral : Agir pessoal e coletivamente com autonomia, responsabilidade, flexibilidade, resiliência e determinação, tomando decisões com base em princípios éticos, democráticos, inclusivos, sustentáveis e solidários.
Articulação com o projeto de vida : Esta competência se articula com o projeto de vida mediante a necessidade dos estudantes de compreenderem que seus projetos de vida devem ser orientados por princípios éticos que não firam a dignidade pessoal, nem a de outras pessoas, além de poderem transformar a realidade pessoal, local e global por meio de projetos de vida comprometidos com a melhoria da qualidade de vida e do meio ambiente.
Habilidades : - Reconhecer as responsabilidades pessoais com as demandas do mundo comum e atribuir sentido ético e sociopolítico ao projeto de vida, comprometendo -se com ações individuais e coletivas voltadas ao bem comum. - Avaliar as implicações éticas das profissões, e sobre elas refletir, compreendendo que é possível contribuir com a sociedade por meio do exercício profissional. - Analisar histórias de vida de agentes responsáveis por transformações sociais, considerando suas particularidades, identidades, diversidades e contextos. - Propor ações de respeito, solidariedade, inclusão e equidade dos múltiplos sujeitos históricos, nos lugares de vivência.
Objetos de conhecimento : - Cidadania participativa. - Trabalho voluntário e atuação das organizações não governamentais ( ONGs ). - Causas e movimentos sociais. - Código de ética das profissões. - Histórias de pessoas que foram/são agentes de transformação social.
Fonte : Elaboração dos autores.

% matched lemmas: acompanhamento, apoiar, autoconfiança, colega, criança, criterioso, cuidado, curiosidade, desejo, equilíbrio, escuta, hábito
\end{textsample}
